\documentclass[aspectratio=169,11pt]{beamer}
\usepackage{../../../shared/theme/esmad_beamer_theme}
\usepackage{amsmath,amssymb}
\usepackage{booktabs}

\title[Robustness and Sensitivity]{Robustness, Placebos, Negative Controls, and Sensitivity Analysis}
\subtitle{Testing the Credibility of Causal Claims}
\date{\today}

\begin{document}

\begin{frame}
\titlepage
\end{frame}

\begin{frame}{Learning Goals}
\begin{itemize}
\item distinguish sampling uncertainty from identification uncertainty;
\item design placebo tests tied to specific failure modes;
\item use negative controls to detect residual confounding and invalid pathways;
\item interpret specification stability across defensible models;
\item evaluate sensitivity to unmeasured confounding;
\item communicate when a causal conclusion is robust, fragile, or not identified.
\end{itemize}
\end{frame}

\begin{frame}{Outline}
\tableofcontents
\end{frame}

\section{Robustness and Identification}

\begin{frame}{Robustness Is Design-Specific}
A robustness check is useful when it probes a plausible violation of the identifying assumptions.

\begin{itemize}
\item DiD: parallel trends and anticipation;
\item IV: exclusion and instrument independence;
\item synthetic control: donor support and pre-treatment fit;
\item observed-confounding adjustment: overlap and omitted confounding.
\end{itemize}

\begin{alertblock}{Principle}
Passing many generic checks is less informative than passing the few checks that directly challenge the design.
\end{alertblock}
\end{frame}

\begin{frame}{Sampling vs Identification Uncertainty}
Standard errors quantify sampling variability under maintained assumptions.

They do not quantify uncertainty about:
\begin{itemize}
\item invalid exclusion restrictions;
\item unmeasured confounding;
\item spillovers;
\item donor contamination;
\item treatment misclassification;
\item non-parallel untreated trends.
\end{itemize}
\end{frame}

\section{Placebo Tests}

\begin{frame}{Placebo Dates}
Move treatment to a false pre-treatment date.

A nonzero placebo effect may indicate:
\begin{itemize}
\item pre-existing trends;
\item anticipation;
\item seasonality;
\item unrelated structural breaks;
\item model misspecification.
\end{itemize}
\end{frame}

\begin{frame}{Placebo Units}
Assign treatment to units known not to have been treated.

Compare their estimated effects with the treated unit's effect.

\begin{block}{Use}
This is especially useful in synthetic-control and panel settings.
\end{block}
\end{frame}

\begin{frame}{Placebo Outcomes}
Estimate an effect on an outcome that treatment should not plausibly affect.

A detected effect can reveal:
\begin{itemize}
\item residual confounding;
\item concurrent shocks;
\item measurement changes;
\item design misspecification.
\end{itemize}
\end{frame}

\section{Negative Controls}

\begin{frame}{Negative-Control Outcomes}
A negative-control outcome shares some confounding structure with the target outcome but should not respond to treatment.

If it appears to respond, the causal design may contain residual bias.
\end{frame}

\begin{frame}{Negative-Control Exposures}
A negative-control exposure resembles the treatment in confounding structure but should not affect the target outcome.

An association can indicate:
\begin{itemize}
\item unmeasured confounding;
\item selection bias;
\item incorrect temporal ordering;
\item data-linkage problems.
\end{itemize}
\end{frame}

\section{Specification Stability}

\begin{frame}{Specification Curves}
A specification curve summarizes estimates across defensible analysis choices:
\begin{itemize}
\item covariate sets;
\item lag lengths;
\item functional forms;
\item comparison groups;
\item trimming thresholds;
\item clustering choices.
\end{itemize}
\end{frame}

\begin{frame}{What Stability Means}
A stable result shows similar sign and magnitude across credible specifications.

Instability can reveal:
\begin{itemize}
\item dependence on one modeling choice;
\item weak support;
\item researcher degrees of freedom;
\item hidden extrapolation.
\end{itemize}

\begin{alertblock}{Caution}
Only substantively defensible specifications belong in the comparison set.
\end{alertblock}
\end{frame}

\section{Sensitivity to Unmeasured Confounding}

\begin{frame}{How Much Confounding Would It Take?}
Sensitivity analysis asks how strong an omitted common cause would need to be to:
\begin{itemize}
\item move the estimate to zero;
\item reverse its sign;
\item push it below a practical decision threshold.
\end{itemize}
\end{frame}

\begin{frame}{Matched-Study Sensitivity}
In matched designs, one can allow matched units to differ in treatment odds because of an unobserved factor.

Inference is recomputed as the allowed difference increases.

\begin{block}{Interpretation}
The more unobserved selection required to overturn the conclusion, the less fragile the result appears under that sensitivity model.
\end{block}
\end{frame}

\begin{frame}{Variance-Based Sensitivity}
Another strategy parameterizes an omitted confounder by the residual variation it explains in:
\begin{itemize}
\item treatment;
\item outcome.
\end{itemize}

Observed covariates can then serve as empirical benchmarks.
\end{frame}

\section{Overlap and Influence}

\begin{frame}{Trimming Sensitivity}
When propensity scores approach zero or one, estimates can depend heavily on a few observations.

Vary trimming thresholds and inspect:
\begin{itemize}
\item effect stability;
\item weight concentration;
\item effective sample size;
\item changes in the target population.
\end{itemize}
\end{frame}

\begin{frame}{Leave-One-Cluster-Out}
Re-estimate after removing one cluster at a time.

This reveals:
\begin{itemize}
\item influential regions or firms;
\item dependence on one treated cluster;
\item leverage hidden by aggregate uncertainty.
\end{itemize}
\end{frame}

\begin{frame}{Leave-One-Donor-Out}
For synthetic control, omit each donor in turn.

If the estimated effect disappears whenever one donor is removed, the result depends strongly on that donor.
\end{frame}

\section{Design-Specific Diagnostics}

\begin{frame}{Difference-in-Differences}
Useful diagnostics include:
\begin{itemize}
\item pre-treatment event-time coefficients;
\item placebo treatment dates;
\item alternative comparison groups;
\item negative-control outcomes;
\item sensitivity to anticipation windows.
\end{itemize}
\end{frame}

\begin{frame}{Instrumental Variables}
Useful diagnostics include:
\begin{itemize}
\item first-stage strength;
\item balance across instrument values;
\item outcomes that should not respond directly to the instrument;
\item subsamples where the treatment channel should be absent;
\item sensitivity to plausible exclusion violations.
\end{itemize}
\end{frame}

\begin{frame}{Synthetic Control}
Useful diagnostics include:
\begin{itemize}
\item pre-treatment fit;
\item placebo treated units;
\item placebo dates;
\item donor-pool restrictions;
\item leave-one-donor-out analysis.
\end{itemize}
\end{frame}

\section{Applied Workflow}

\begin{frame}{Example: Regional Pricing Policy}
Suppose the estimated effect is a 5 percent reduction in demand.

A useful robustness program asks:
\begin{itemize}
\item Is there an effect at false pre-treatment dates?
\item Does an unrelated outcome move too?
\item Does the effect survive alternative comparison regions?
\item Does it depend on one region?
\item Does trimming poor-overlap regions matter?
\item How strong would omitted confounding need to be to remove the effect?
\end{itemize}
\end{frame}

\begin{frame}{A Credible Robustness Workflow}
\begin{enumerate}
\item List the identifying assumptions.
\item Map each assumption to plausible failure modes.
\item Choose targeted falsification tests.
\item Re-estimate across defensible specifications.
\item Quantify influence of key units, clusters, or donors.
\item Run sensitivity analysis for unmeasured confounding.
\item Separate estimator sensitivity from estimand changes.
\item Report failed checks prominently.
\item Revise the causal claim when diagnostics weaken identification.
\end{enumerate}
\end{frame}

\begin{frame}{Common Failure Modes}
\begin{itemize}
\item reporting only checks that pass;
\item running irrelevant checks for appearance;
\item treating a null placebo as proof of identification;
\item hiding dependence on one specification;
\item ignoring failed negative controls;
\item calling different estimands “robustness checks” without noting the target changed.
\end{itemize}
\end{frame}

\begin{frame}{Synthesis: Robustness Is About Credibility}
\begin{itemize}
\item Robustness analysis should target plausible violations of identifying assumptions.
\item Placebos ask whether effects appear where they should not.
\item Negative controls probe residual confounding and invalid pathways.
\item Specification curves reveal dependence on defensible modeling choices.
\item Sensitivity analysis quantifies how strong omitted confounding must be to alter conclusions.
\item Influence diagnostics expose dependence on individual units, clusters, or donors.
\item Failed diagnostics should weaken the causal claim, not be buried.
\end{itemize}
\end{frame}

\begin{frame}{Selected References}
\small
\begin{itemize}
\item Rosenbaum, P. R. (2002). \textit{Observational Studies}, 2nd ed. Springer.
\item Cinelli, C., Ferwerda, J., and Hazlett, C. (2020). Sensitivity Analysis Tools for OLS. \textit{Journal of Statistical Software}, 96(1), 1--28.
\item Simonsohn, U., Simmons, J. P., and Nelson, L. D. (2020). Specification Curve Analysis. \textit{Nature Human Behaviour}, 4, 1208--1214.
\item Lipsitch, M., Tchetgen Tchetgen, E., and Cohen, T. (2010). Negative Controls: A Tool for Detecting Confounding and Bias in Observational Studies. \textit{Epidemiology}, 21(3), 383--388.
\item Abadie, A. (2021). Using Synthetic Controls: Feasibility, Data Requirements, and Methodological Aspects. \textit{Journal of Economic Literature}, 59(2), 391--425.
\end{itemize}
\end{frame}

\contactslide

\end{document}
